\section{Lexical structure}
\label{sec:lexical}

\subsection{Source text}

A program is a sequence of bytes, conventionally UTF-8. The lexer works
on bytes; only ASCII is significant to the grammar. Both \ic{"\bs n"} (LF) and \ic{"\bs r\bs n"} (CRLF) end a line. Positions reported
in diagnostics are 1-based line and column, the column counted in bytes.

\subsection{Comments}

\begin{lstlisting}
// a line comment runs to the end of the line
/* a block comment
   does not nest: the first */ closes it */
\end{lstlisting}

\noindent Comments are whitespace: they separate tokens and nothing more.

\subsection{Identifiers and keywords}

\ic{isalpha}(c) is true for \ic{A}--\ic{Z}, \ic{a}--\ic{z} and
\ic{\_}; \ic{isdigit}(c) for \ic{0}--\ic{9}. An identifier is

\begin{align*}
\ty{identifier} &::= \ty{isalpha}~\{~\ty{isalnum}~\}
\end{align*}

\noindent so \ic{\_} is a word character: \ic{is\_digit} is a single
identifier, not three tokens. The prefix \ic{y\_} is \emph{reserved}:
identifiers beginning with \ic{y\_} are used by the compiler's generated
code and must not be written by hand.

The keywords are:
\begin{center}
\ic{let} \ic{const} \ic{def} \ic{if} \ic{else} \ic{while} \ic{for}
\ic{foreach} \ic{in} \ic{return} \ic{true} \ic{false}
\end{center}

\subsection{Literals}
\label{sec:literals}

\paragraph{Integers.} Decimal, binary (\ic{0b1010}) and hexadecimal
(\ic{0x1F}) radix literals are supported. Any of them may group digits
with \ic{\_}, which is ignored: a radix literal (\ic{0x} / \ic{0b},
either case) in any position after the prefix (\ic{0xFF\_FF},
\ic{0b1010\_1010}), a decimal one anywhere after its first digit
(\ic{1\_000\_000}; a leading \ic{\_} starts an identifier instead).
A prefix needs a digit of its radix: \ic{0x}, \ic{0b\_} and \ic{0b2}
are errors (\ic{binary literal '0b2' has no digits}), one for the whole
word.
There is no octal form: \ic{017} is decimal seventeen. All integer
literals are \ty{long} (64-bit). A radix literal spells the 64 bits
themselves, so a mask may set the top bit:
\ic{0xFFFF\_FFFF\_FFFF\_FFFF} is \ic{-1} and
\ic{0x8000\_0000\_0000\_0000} the smallest \ty{long}; digits past 64
bits wrap away (modulo $2^{64}$).

\paragraph{Floats.} A digit sequence with a decimal point and/or an
exponent: \ic{1.5}, \ic{1.5e+3}, \ic{2.0e-1}. \ic{\_} may group digits
after the first digit of the whole part, of the fraction and of the
exponent: \ic{1\_000.5}, \ic{0.000\_001}, \ic{1e1\_0} (\ic{1.\_5} and
\ic{1e\_5} are parse errors). A dot followed by a non-digit is
\emph{not} a fractional part --- \ic{a..b} lexes as \ic{a}, \ic{..},
\ic{b}. Float literals are IEEE-754 binary64 (\ty{float}).

\paragraph{Strings.} Double-quoted byte strings:

\begin{center}
\begin{tabular}{ll}
\toprule
escape & meaning \\
\midrule
\ic{\bs n} \ic{\bs t} \ic{\bs r} & LF, HT, CR \\
\ic{"\bs "} \ic{"\bs ""} \ic{"\bs '"} & backslash, quote \\
\ic{\bs xHH} & byte with hex value \ic{HH} \\
\ic{\bs c} (any other \ty{c}) & \ty{c}; the backslash is dropped \\
\bottomrule
\end{tabular}
\end{center}

\noindent An escape is read as a single byte, and an unrecognised
\ic{\bs c} simply drops the backslash (\ic{"q\bs q"} is \ic{"qq"}), so a
NUL is spelled \ic{"\bs x00"} --- which is not \ic{"\bs0"} (that is the
character \ic{0}). A \ic{str} is a length-prefixed byte string, so an
escaped NUL inside one is ordinary data and \ic{len} counts it (use
\ic{[byte]} arrays, Section~\ref{sec:builtin-bytes}, for bulk binary
data).

\paragraph{Byte strings.} \ic{b"..."} and \ic{b'...'} are \emph{byte
string literals}: a \ic{[byte]} value spelled like a Python \ic{bytes}
literal. The escapes are the string escapes above, so binary data writes
naturally: \ic{b"\bs x7fELF"} is the four-byte ELF magic. Either quote
spells the same type; single quotes exist only for byte strings.

\paragraph{Format strings.} \ic{f"..."} is a string literal with
\emph{fields}, as in Python:
\begin{center}
\ic{f"\{name\} is \{age\} (\{share:.1\%\})"}
\end{center}
\noindent A field \ic{\{e\}} becomes \ic{tostring(e)}, and
\ic{\{e:spec\}} becomes \ic{e} formatted as \ty{spec} says
(Section~\ref{sec:format-spec}). The text around the fields takes the
string escapes, and \ic{\{\{} and \ic{\}\}} (or \ic{\bs\{} and
\ic{\bs\}}) are literal braces. An f-string is sugar for concatenation:
\ic{f"a\{x\}b"} means \ic{"a" + tostring(x)} \ic{+ "b"}, a \ty{str} like
any other.

A field holds one expression --- calls, string literals, record literals
and nested f-strings included --- and may span lines; an error inside it
points at its own line and column. The expression ends at the first
\ic{\}} or \ic{:} outside brackets, except a \ic{:} that answers a
\ic{?}: \ic{\{ok ? a : b\}} works, while a cons has to be parenthesised,
\ic{\{(x :: xs)\}}. A \ic{?} that is part of a \ic{?.} (Section~
\ref{sec:opt-propagation}) answers nothing, so the \ic{:} after it is
the spec's own: \ic{f"\{o?.port:x\}"\} is \ic{o?.port} as \ic{x}. The spec
is plain text up to the field's \ic{\}}; it cannot hold another field.

\paragraph{Booleans and none.} \ic{true}, \ic{false}. The empty option
is spelled \ic{none}, a predeclared name rather than a keyword
(Section~\ref{sec:type-options}). There is no \ic{null}: the word is an
ordinary (undefined) identifier.

\subsection{Operators and punctuation}

\begin{center}
\begin{tabular}{lll}
\toprule
kind & tokens \\
\midrule
% `<{}<`: the {} keeps T1's << and >> ligatures (guillemets) out.
three- and four-character &
  \ic{>{}>{}>} \ic{<{}<=} \ic{>{}>=} \ic{>{}>{}>=} \\
two-character &
  \ic{==} \ic{!=} \ic{<=} \ic{>=} \ic{\&\&} \ic{||} \ic{::} \ic{..}
  \ic{->} \ic{=>} \ic{++} \ic{--} \ic{<{}<} \ic{>{}>} \\
 &
  \ic{+=} \ic{-=} \ic{*=} \ic{/=} \ic{\%=}
  \ic{\&=} \ic{|=} \ic{\textasciicircum{}=} \\
single-character &
  \ic{+} \ic{-} \ic{*} \ic{/} \ic{\%} \ic{<} \ic{>} \ic{=} \ic{!}
  \ic{\&} \ic{.} \ic{?} \ic{|} \ic{\textasciicircum{}} \ic{\textasciitilde{}} \\
punctuation &
  \ic{(} \ic{)} \ic{\{} \ic{\}} \ic{[} \ic{]} \ic{,} \ic{;} \ic{:} \\
\bottomrule
\end{tabular}
\end{center}

\paragraph{Maximal munch.} The lexer always takes the longest operator
that matches, exactly like C:
\begin{itemize}[nosep]
\item \ic{x+++y} is \ic{x ++ + y}, not \ic{x + ++y};
\item \ic{x--y} is \ic{x-- y} (a parse error), not \ic{x - -y};
      write \ic{x - -y} when negation is meant;
\item \ic{/=} is one token; \ic{//} starts a comment, and
      \ic{/ == y} must be written with a space to divide;
\item \ic{x>{}>{}>=1} is the compound \ic{x >{}>{}>= 1};
\end{itemize}

\paragraph{Contextual keywords.} \ic{none}, \ic{match}, \ic{type} and
\ic{void} are ordinary identifiers to the lexer: \ic{none} is a
predeclared name (Section~\ref{sec:type-options}), \ic{match} is the
match-expression keyword only when an expression and \ic{\{} follow it,
and \ic{type} starts a type alias only when followed by \ic{name =}. Each
stays usable as an ordinary name elsewhere.
